\scanpage{636}
\mspage{24}

\unsure{donc} établir Der 4, Der 5 : existence des $u_*^W$, et
calcul des fibres de $u_*^W(X)$. Quant à Der 6,
il est acquis automatiquement, du fait
que chaque $(M(I),W(I))$ s'insère dans une
structure de fibrations $(W(I),F(I))$. (Oublié
dans la démonstration de 2.4.) L'idée naturelle
est d'utiliser à présent la structure de
triple

\[
 \left(W(I),C(I),\underset{\substack{\|\,\mathrm{def}\\ TC(I)_{*}}}{F[I]}\right),\tag*{(2.5.1)}
\]

mais la difficulté est de prouver la \unsure{factorisabilité|factorisabili-}
pour les couples

\[
 C(I)\underset{\unsure{<|?}}{\leftrightarrow}TF[I],\quad
 \underset{\substack{\|\,\mathrm{def}\\ C(I)\cap W(I)}}{TC(I)}\longleftrightarrow F[I],\tag*{(2.5.2)}
\]

faute de pouvoir montrer, dans les conditions
générales de 2.3 seulement (triple de Quillen
spécial) que $C(I)$ et $TC(I)$ sont également
spéciaux. (Cf. \unsure{question} 1.11.) Mais
moyennant des conditions convenables sur
$M$, on arrive à s'en tirer si $C=\underline{\mathrm{mono}}$.
L'énoncé pertinent, un peu technique, se
trouve dans \uline{XV} cor 7.4 (p 98--99), que je
vais rappeler ici.
% archive sheet number: 635
